% ============================================================
%  ch17  附录 A §A.1 层论补充：等变反常层
%  原文行 3405–3537
% ============================================================

\appendix

\section{层论补充}\label{App}

\subsection{等变反常层范畴}

\subsubsection{基本性质}
回忆对 $k$ 上的概形 $X$，有可构造 $\ell$-进复形的有界导出范畴 $D_c^b(X)$，其中反常层范畴 $\on{P}(X)$ 作为全 Abel 子范畴。若 $X$ 带有线性代数群 $K$ 的作用（作用映射记作 $\on{act}$），则可以定义 $X$ 上 $K$-等变反常层的 Abel 范畴 $\on{P}_K(X)$，即：$\on{P}_K(X)$ 的对象是 $X$ 上的反常层 $\mF$ 连同一个沿两个映射
\[\on{act},\ \pr_2:K\times X\longrightarrow X\]
的同构 $\theta$，满足 (i) $e^*\theta=\id$，其中 $e:X\to K\times X$ 由 $K$ 的单位元给出；(ii) $K\times K\times X$ 上自然的余圈条件。

\begin{rmk}
还有更复杂的概念：$X$ 上（可构造）$\ell$-进层的（有界）等变导出范畴 $\on{D}_{K,c}^b(X)$，它不能像上面那样简单地定义。或者如 \cite{BeLu} 定义，或者把它看作代数叠 $[X/K]$ 上 $\ell$-进层的导出范畴（例如参见 \cite{LZ}）。本书只在非常形式的层面偶尔用到这个概念，即这些三角范畴之间（在不同概形上）存在良好的 Grothendieck 六函子形式体系。\zy{该六函子形式体系的一般框架见 nLab 词条 six operations（\url{https://ncatlab.org/nlab/show/six+operations}，2026-09-30 修订）。[来源：nLab six operations，2026-09-30。]}
\end{rmk}

按定义有遗忘函子 $\on{P}_K(X)\to \on{P}(X)$。

\begin{lem}\label{A:forgetful}
若 $K$ 连通，则该遗忘函子完全忠实，其本质像由那些满足 $\on{act}^*\mF\cong\pr_2^*\mF$ 的反常层 $\mF$ 构成。
\end{lem}

本引理的意义在于：$K$-等变性是 $X$ 上反常层的一个\emph{性质}，而非附加结构。注意该陈述对不连通群不成立：若 $K$ 有限，则 $\on{P}_K(\on{pt})$ 等价于 $K$ 的有限维表示范畴。该陈述对等变导出范畴也不成立。

\begin{proof}
简述第一个断言的证明，第二个留作练习。该函子显然忠实。设 $(\mF_1,\theta_1),(\mF_2,\theta_2)$ 是两个 $K$-等变反常层，$\varphi: \mF_1\to \mF_2$ 是态射，要证它与 $\theta_i$ 自动相容，即
\begin{equation}\label{A:comp}
\theta_2^{-1}\circ\on{act}^*(\varphi)\circ\theta_1=\pr_2^*(\varphi).
\end{equation}
注意因 $K$ 连通，由 \cite[Proposition 4.2.5]{BBD}，$\pr^*[\dim K]: \on{P}(X)\to \on{P}(K\times X)$ 完全忠实。于是要验证 \eqref{A:comp}，只需证 $e^*(\theta_2^{-1}\circ\on{act}^*(\varphi)\circ\theta_1)=e^*(\pr_2^*(\varphi))$，而这按定义成立。
\end{proof}

我们将用到等变范畴的下述两条性质。

\begin{lem}\label{A:unip}
设 $K_1\subset K$ 是闭正规子群。
\begin{enumerate}
\item 若 $K_1$ 在 $X$ 上的作用自由，且 $[X/K_1]$ 可由代数空间 $\bar X$ 表示，则沿 $q:X\to \bar X$ 的拉回诱导范畴等价
\begin{equation}\label{equivariance free action}
q^*[\dim K_1]:\on{P}_{K/K_1}(\bar X)\simeq \on{P}_K(X).
\end{equation}

\item 设 $K_1$ 连通。若 $K_1$ 在 $X$ 上的作用平凡，则遗忘函子
\begin{equation}\label{unipotence}
\on{P}_K(X)\to \on{P}_{K/K_1}(X)
\end{equation}
是范畴等价。
\end{enumerate}
\end{lem}

\begin{proof}
由反常层的 \'etale 下降，可假设 $X=\bar X\times K_1$，于是 (1) 容易推出。(2) 由引理 \ref{A:forgetful} 推出。
\end{proof}

\begin{rmk}\label{equivariant descent}
若把 $\on{P}_{K/K_1}$ 换成等变导出范畴，等价 \eqref{equivariance free action} 仍成立，但证明更不平凡（见 \cite[\S\ 2.6.2]{BeLu}）。
\end{rmk}

\subsubsection{扭曲外积}\label{A:twist product}
设 $p:E\to X$ 是 $K$-主丛，$Y$ 是带 $K$ 作用的概形。设 $\mF\in\on{P}(X)$、$\mG\in \on{P}_K(Y)$，则可如下构造 $X\tilde\times Y$ 上的反常层，即扭曲外积 $\mF\tilde\boxtimes\mG$。$E\times Y$ 上的反常层 $p^*\mF[\dim K]\boxtimes \mG$ 按构造是 $K$-等变的。于是由引理 \ref{A:unip} (i)，它下降为 $X\tilde\times Y$ 上唯一（在唯一同构意义下）的反常层，即 $\mF\tilde\boxtimes\mG$。注意由注 \ref{equivariant descent}，外扭曲积其实可对 $\ell$-进复形定义。

\subsubsection{等变上同调}\label{A:char class}
设 $K$ 是 $k$ 上的线性代数群，$R_{K}=\on{H}^*(\bB K)$ 是 $K$ 的分类叠的上同调。它是有限生成的分次增广交换 $\Ql$-代数，同构于多项式代数。具体地，它可实现如下。设 $\{E_n\to B_n\}$ 是一列 $K$-主丛，逼近 $K$ 的分类空间。例如把 $K$ 嵌入某个 $\GL_r$ 使 $\GL_r/K$ 拟仿射；对大的 $n$，令 $E_n:=S_{n,r}$ 为 Stiefel 簇，即 $\Gr(r,n)$ 上的通用 $\GL_r$-主丛。则 $B_n:=E_n/K$ 可由概形表示，$B_n\subset B_{n+1}$ 是闭嵌入，且 $\on{H}^*(\underrightarrow\lim_n B_n)= \on{H}^*(\bB K)$。例如若 $K=T$ 是环面，则有典范同构 $\Spec R_T\cong \Lie T$。若 $K$ 幂幺，则 $R_K=\Ql$。

设 $\mF\in \on{P}_K(X)$（或更一般地 $\mF\in\on{D}^b_{K,c}(X)$），则可谈 $K$-等变上同调 $\on{H}_K^*(X,\mF)$。它或者定义为 $\mF$ 作为商叠 $[X/K]$ 上层的上同调，或者更具体地
\[\on{H}_K^*(X,\mF):=\on{H}^*(\underrightarrow\lim_n (B_n\tilde\times X), \Ql\tilde\boxtimes \mF).\]
由构造，$\on{H}^*_K(X,\mF)$ 是 $\on{H}^*(\underrightarrow\lim_n B_n)=R_{K}$ 上的模。

由定义可知：若 $K_1\subset K$ 自由作用于 $X$，$\bar X$ 为如引理 \ref{A:unip} 中的商，则
\begin{equation}\label{coh free action}
\on{H}^*_K(X,q^*\mF)=\on{H}^*_{K/K_1}(\bar X,\mF).
\end{equation}
另一方面，若 $K_1$ 幂幺且平凡作用于 $X$，则
\begin{equation}\label{coh unip}
\on{H}^*_K(X,\mF)=\on{H}^*_{K/K_1}(X,\mF).
\end{equation}

现在设 $K$ 连通，$\mF\in\on{P}_K(X)$。与投影 $B_n\tilde\times X\to B_n$ 相伴的 Leray 谱序列诱导一个谱序列，其 $E_2$-项为 $\on{H}^*(\bB K)\otimes \on{H}^*(X,\mF)$，收敛到 $\on{H}^*_K(X,\mF)$。仿照 \cite{GKM}，称 $\mF$ \textbf{等变形式}（equivariantly formal），若该谱序列在 $E_2$-项退化。

\begin{thm}
若 $X$ proper，则 $X$ 的相交上同调层（看作 $X$ 上 $K$-等变反常层）等变形式。
\end{thm}

证明依赖「权」的理论，见 \cite[Theorem 14.1]{GKM}。事实上同一论证适用于几何来源的半单纯反常层。

对等变形式的复形 $\mF$，有非典范的同构 $\on{H}^*_K(X,\mF)\cong R_{K}\otimes\on{H}^*(X,\mF)$。特别地此时 $\on{H}^*_K(X,\mF)$ 是秩等于 $\dim \on{H}^*(X,\mF)$ 的自由 $R_{K}$-模。典范地有同构
\begin{equation}\label{de-equiv}
\on{H}^*(X,\mF)\cong \on{H}^*_K(X,\mF)\otimes_{R_{K}}\Ql,
\end{equation}
其中 $R_{K}\to\Ql$ 是增广映射。此外有等变 K\"unneth 公式：设 $X$、$Y$ 都 proper，$\mF\in\on{P}_K(X)$、$\mG\in\on{P}_K(Y)$ 是半单纯反常层，则有典范同构
\begin{equation}\label{equiv Kunneth}
\on{H}_K^*(X,\mF)\otimes_{R_{K}}\on{H}_K^*(Y,\mG)\cong \on{H}^*_K(X\times Y,\mF\boxtimes\mG).
\end{equation}
这由通常的 K\"unneth 公式与谱序列的退化推出。

再叙述等变上同调的局部化定理（拓扑版本见 \cite[(6.2)]{GKM}）。首先对 $\mF\in \on{D}_{T,c}^b(X)$ 可定义紧支撑 $T$-等变上同调为
\[\on{H}_{T,c}^*(X,\mF):=\on{H}_c^*(\underrightarrow\lim_n (B_n\tilde\times X), \Ql\tilde\boxtimes \mF).\]

\begin{thm}\label{equiv localization}
设 $K=T$ 是作用于 $X$ 的环面，$i:X^0\subset X$ 是 $T$-不动点子概形。则对 $\mF\in\on{D}_{T,c}^b(X)$，自然映射
\[\on{H}^*_T(X^0,i^!\mF)\to \on{H}^*_T(X,\mF),\quad \on{H}^*_{T,c}(X,\mF)\to \on{H}^*_{T,c}(X^0,i^*\mF)\]
在局部化到 $R_{T}$ 的通用点后成为同构。
\end{thm}

\begin{proof}
简述第一个同构的证明，第二个由 Verdier 对偶推出。设 $j: X-X^0\to X$ 是开补集。由 distinguished triangle $i_!i^!\mF\to \mF\to j_*j^*\mF\to$，化归到证明：若 $X^0=\emptyset$，则 $\on{H}^*_T(X,\mF)$ 是挠 $R_T$-模。可取 $X$ 的有限覆盖 $X=\cup X_i$，使每个 $X_i$ 是 $T$-不变的，且存在非平凡子环面 $T_i\subset T$ 自由作用于 $X_i$。于是由 \eqref{coh free action}，$\on{H}^*_T(X_i,\mF|_{X_i})$ 是挠 $R_T$-模。从而由谱序列论证，$\on{H}^*_T(X,\mF)$ 也是挠 $R_T$-模。
\end{proof}

\subsubsection{作用于归纳概形的 pro-代数群}\label{A:ind}
设 $X$ 是归纳概形。可把 $X$ 上的反常层范畴定义为 $\on{P}(X)=\underrightarrow\lim\on{P}(X_i)$，其中 $X=\underrightarrow\lim X_i$ 是把 $X$ 表为 $k$ 上有限型闭子概形的表示。这与表示的选取无关。

现设 $X$ 带有 pro-代数群 $K$ 的作用。假设
\[(S)\quad\quad \left.\begin{array}{c}\mbox{每个几何点的稳定化子只有有限多个连通分支。}\end{array}\right.\]
则可以如下定义 $X$ 上 $K$-等变反常层的范畴。

首先存在表示 $X=\underrightarrow\lim X_i$，使每个 $X_i$ 是 $k$ 上有限型的 $K$-稳定闭子概形，且 $K$ 在 $X_i$ 上的作用通过某个代数商 $K_i$ 分解。可取 $K_i$ 为 $K$ 模去连通正规子群所得的代数商。此外，若 $X_i\subset X_j$ 是闭子概形，可安排 $K_i$ 是 $K_j$ 的商。

然后定义 $\on{P}_K(X_i):=\on{P}_{K_i}(X_i)$。由引理 \ref{A:unip} (ii) 可知 $\on{P}_K(X)$ 在典范等价意义下与 $K_i$ 的选取无关。因此有完全忠实函子
\[\on{P}_K(X_i)=\on{P}_{K_i}(X_i)\simeq \on{P}_{K_j}(X_i)\to \on{P}_{K_j}(X_j)=\on{P}_K(X_j).\]
最后定义
\[\on{P}_K(X)=\underrightarrow\lim \on{P}_K(X_i).\]

现设 $E\to X$ 是有限型归纳概形上的 $K$-主丛，$Y$ 是满足 (S) 的 $K$-归纳概形。设 $\mF\in\on{P}(X)$、$\mG\in\on{P}_K(Y)$，可类似地把扭曲积 $\mF\tilde\boxtimes\mG$ 定义为 $\on{P}(X\tilde\times Y)$ 的对象。即：$\mF$ 支在某个 $X_i$ 上，$\mG$ 支在某个 $Y_j$ 上，$K$ 在 $Y_j$ 上的作用通过 $K_j$ 分解。则 $E_{ij}:=E|_{X_i}\times^KK_i\to X_i$ 是 $K_i$-主丛，且
\[X_i\tilde\times Y_j= E|_{X_i}\times^KY_j= E_{ij}\times^{K_j}Y_j\]
是 $X\tilde\times Y$ 的闭子概形。于是 \S\ref{A:twist product} 的构造给出 $\on{P}(X_i\tilde\times Y_j)\subset \on{P}(X\tilde\times Y)$ 中的对象 $\mF\tilde\boxtimes\mG$。

最后，设 $K^u$ 是 $K$ 的 pro-幂幺根，且 $K/K^u$ 代数。设 $X$ 是满足 (S) 的 $K$-归纳概形。则对 $\mF\in \on{P}_K(X)$ 有定义良好的 $K$-等变上同调 $H_K^*(X,\mF)$，它是 $R_{K}=\on{H}^*(\bB (K/K^u))$ 上的模。\S\ref{A:char class} 中关于等变上同调的全部性质在此情形下仍成立。
